\documentclass[11pt]{article}

% --- Packages (The "plugins" of LaTeX) ---
\usepackage[utf8]{inputenc} % Character encoding
\usepackage{geometry}       % Page margins
 \geometry{margin=1in}
\usepackage{graphicx}       % For inserting images
\usepackage{amsmath}        % For math equations
\usepackage{listings}
\usepackage{xcolor} % Required for syntax highlighting colors
\usepackage{booktabs} % For professional lines
\usepackage{siunitx}  % For decimal alignment

% Define colors for Python syntax
\definecolor{codegreen}{rgb}{0,0.6,0}
\definecolor{codegray}{rgb}{0.5,0.5,0.5}
\definecolor{codepurple}{rgb}{0.58,0,0.82}
\definecolor{backcolour}{rgb}{0.95,0.95,0.92}
\definecolor{termback}{rgb}{0.1,0.1,0.1}
\definecolor{termtext}{rgb}{0.9,0.9,0.9}
\definecolor{lightgray}{rgb}{0.9,0.9,0.9}

\lstdefinestyle{mystyle}{
    backgroundcolor=\color{backcolour},   
    commentstyle=\color{codegreen},
    keywordstyle=\color{magenta},
    numberstyle=\tiny\color{codegray},
    stringstyle=\color{codepurple},
    basicstyle=\ttfamily\footnotesize,
    breakatwhitespace=false,         
    breaklines=true,                 
    captionpos=b,                    
    keepspaces=true,                 
    numbers=left,                    
    numbersep=5pt,                  
    showspaces=false,                
    showstringspaces=false,
    showtabs=false,                  
    tabsize=2
}

\lstdefinestyle{terminal}{
    backgroundcolor=\color{white},   % White background
    basicstyle=\ttfamily\small\color{black}, % Black text
    frame=single,                    % Simple border
    rulecolor=\color{lightgray},     % Light gray border color
    breaklines=true,                 % Wrap long lines
    numbers=none,                    % No line numbers for output
    captionpos=b,                    % Caption at the bottom
    xleftmargin=10pt,                
    xrightmargin=10pt,
    framesep=5pt
}

\lstset{style=mystyle}

% --- Document Information ---
\title{DE 414: Practical 2 Report}
\author{Amy McDermott (26911264)}
\date{\today}

\begin{document}

\maketitle

\section*{Introduction}
This report documents the implementation of a logistic regression model using Python. The model is trained on the MNIST dataset with the aim to 
correctly classify digits. This report includes the code used for training the model, as well as the results obtained.
-describe question posed
- strategy to solve the
- software designed to implement solution
- analysis of results that were obtained

\section*{Question 1: Softmax Function}
The softmax function is a common activation function used to convert a vector of raw scores (logits) into a probability distribution, with values ranging from 0 to 1 that sum to 1. The function is defined as follows: 
\begin{equation}
\sigma(\mathbf{z})_i = \frac{e^{z_i}}{\sum_{j=1}^{K} e^{z_j}} \ \ \ \text{for } i = 1, \dots, K
\end{equation}
where $\mathbf{z}$ is the logit of class $i$ and $K$ is the total number of classes.
The output of the softmax function yields interpretable probability scores for each class.

The implementation of this function in Python, computes softmax for an N-dimensional NumPy array \texttt{x} along a specified axis (defaulting to the last dimension). It uses vectorized operations for efficiency on large datasets.
To ensure numerical stability, the maximum value along the axis is subtracted from the input before exponentiation to prevent overflow/underflow issues. This shifts values without altering the relative probabilities.
The function then computes the exponentials of the shifted inputs and normalizes them by dividing by the sum of exponentials along the specified axis, resulting in a probability distribution. The 
listing \ref{lst:softmax} below shows the implementation of the softmax function in Python.

\begin{lstlisting}[
    language=Python, 
    caption={Softmax function implementation},
    label={lst:softmax},
    breaklines=true,
    basicstyle=\ttfamily\small,
    frame=single
]
shift_X = X - np.max(X, axis = axis, keepdims = True) # Subtract the maximum for numerical stability
exps = np.exp(shift_X) # Compute the exponentials
return exps / np.sum(exps, axis = axis, keepdims = True) # Compute softmax
\end{lstlisting}

The vectorized implementation is a design choice that leverages Numpy's operations for fast computation without loops. This is suitable for large datasets.

\section*{Question 2: Log-likelihood Function}
A function is implemented to compute the log-likelihood between one-hot encoded target vectors $\mathbf{y}$, and predictions $\mathbf{\hat{y}}$. This function computes how well predicted probabilities match the true one-hot encoded labels 
in multiclass classification. The training set log-likelihood is defined as:
\begin{equation}
\ell(\mathbf{W}) = \sum_{n=1}^{N} \sum_{k=1}^{K} y_{k}^{(n)} \ln \left( \hat{y}_{k}^{(n)} \right)
\end{equation}
The implementation in Python computes the log-likelihood for each sample in a batch and ensures numerical stability by avoiding undefined logarithms. This is achieved by clipping predicted probabilities  
\texttt{Y\_hat} to [epsilon, 1] to prevent log(0), which is undefined and causes errors. The function then performs an element-wise logarithm operation to the clipped predictions, yielding log-probabilities.
for each class per sample. Finally, the log-likelihood is computed by summing the element-wise product of the one-hot encoded targets and the log-probabilities across all samples and classes. This returns 
an array of log-likelihoods for each sample in the batch.
 The implementation is shown in listing \ref{lst:loglikelihood}.
 \begin{lstlisting}[
    language=Python, 
    caption={Log-likelihood function implementation},
    label={lst:loglikelihood},
    breaklines=true,
    basicstyle=\ttfamily\small,
    frame=single
]
epsilon = 1e-15 # Perturb Y_hat to prevent log(0) which is undefined
Y_hat = np.clip(Y_hat, epsilon, 1)
log_probs = np.log(Y_hat) # Compute the log of the predicted probabilities
LL = np.sum(Y * log_probs, axis = 1) # Element-wise multiplication of Y and log_probs, then sum across classes for each sample
return LL # Log likelihoods per sample
\end{lstlisting}

\section*{Question 3: Extract the minibatch from the training set}
Mini-batch training is a common technique in machine learning to balance computational efficiency with convergence stability. 
Mini-batch training uses a subset of the training data to compute the gradient and update model parameters, leading to more frequent updates.
A function is implemented to extract a minibatch from the dataset using a default batch size of 60. 
The variable $\mathbf{i}$ represents the specified index of the minibatch to be extracted. The software design thus computes the start and end index for a 
specified minibatch index based on the batch size, and slices the data matrix $\mathbf{X}$ and target vector $\mathbf{Y}$ to return the corresponding minibatch.
Listings \ref{lst:minibatch} shows the implementation of the minibatch extraction function in Python.

 \begin{lstlisting}[
    language=Python, 
    caption={Minibatch extraction function implementation},
    label={lst:minibatch},
    breaklines=true,
    basicstyle=\ttfamily\small,
    frame=single
]
start = i * batch_size # Starting index of the batch
end = start + batch_size # Ending index of the batch
X_mini = X[start:end] # Slicing along first axis (rows) # Slice the data matrix X and vector Y to get the minibatch
Y_mini = Y[start:end]
return X_mini, Y_mini
\end{lstlisting}

\section*{Question 4: Multinomial logistic regression model}
\subsection*{4.a: Forward pass}
A forward pass function is implemented to compute the predicted probabilities $\mathbf{Y\_hat}$ for a given input $\mathbf{X}$ and model parameters $\mathbf{W}$.
Engineering design includes adding a bias term to the input data by appending a column of ones to the left of the input matrix $\mathbf{X}$. This allows the model to learn an intercept term for each class.
The linear layer is computed by performing a matrix multiplication between the input data with bias and the weight matrix $\mathbf{W}$, resulting in the logits for each class.
The softmax activation function is then applied to the logits to convert them into predicted probabilities $\mathbf{Y\_hat}$ for each class. The implementation of the forward pass function in Python is shown in listing \ref{lst:forwardpass}.
 \begin{lstlisting}[
    language=Python, 
    caption={Forward pass function implementation},
    label={lst:forwardpass},
    breaklines=true,
    basicstyle=\ttfamily\small,
    frame=single
]
# Add bias term to input data (column of 1s to X matrix)
ones = np.ones((X.shape[0], 1)) # Create a column of ones with the same number of rows as X
X_bias = np.hstack([ones, X]) # Append the column of ones to the left of X
logits = X_bias @ self.weights
Y_hat = softmax(logits, axis = 1)
return Y_hat
\end{lstlisting}

\subsection*{4.b: Gradient descent optimisation}
In the class function train, a gradient descent optimisation step is implemented. The model weights $\mathbf{W}$ are updated iteratively based on the computed gradients of the log-likelihood loss with respect to the weights.
The training loop iterates for a specified number of epochs, and in each epoch, the training data is processed in mini-batches. For each mini-batch, the forward pass is computed to obtain predicted probabilities $\mathbf{Y\_hat}$.
The training set log-likelihood for the nth feature xn with respect to the weight vector wk associated with the kth node is computed as: 
\begin{equation}
\frac{\partial \ell^{(n)}}{\partial \mathbf{w}_k} = \left( y_{k}^{(n)} - \hat{y}_{k}^{(n)} \right) \mathbf{x}^{(n)}
\end{equation}
Thus the gradient for all weights $\mathbf{X}$ (average log-likelihood) is computed as: \begin{equation}
\nabla_{\mathbf{W}} \ell = \frac{1}{N} \mathbf{X}^T \left( \mathbf{Y} - \mathbf{\hat{Y}} \right)
\end{equation}
The weights are then updated as follows: 
\begin{equation}
\mathbf{W}_{\text{new}} = \mathbf{W}_{\text{old}} + \eta \nabla_{\mathbf{W}} \ell
\end{equation}
where $\eta$ is the learning rate.

\begin{lstlisting}[
    language=Python, 
    caption={Gradient descent function implementation},
    label={lst:gradientdescent},
    breaklines=true,
    basicstyle=\ttfamily\small,
    frame=single
]
num_samples = train_X.shape[0]
num_batches = num_samples // batch_size 

for epoch in range(epochs):
    for i in range(num_batches):
        x_batch, y_batch = get_mini_batch(train_X, train_Y, i, batch_size)

        y_hat = self.forward(x_batch) # Forward pass to get predictions for the batch

        ones = np.ones((x_batch.shape[0], 1)) # Create a column of ones with the same number of rows as x_batch
        x_batch_bias = np.hstack([ones, x_batch]) # Append the column of ones to the left of x_batch

        # Calculate the gradient of the log-likelihood wrt. to the weights
        grad = (x_batch_bias.T @ (y_batch - y_hat)) / batch_size # Average over the batch

        self.weights += learning_rate * grad
\end{lstlisting}

\subsection*{4.c: Train model}
An instance of the MultinomialLogisticRegression class is created, and the model is trained on the training data by calling the train method for all batches in the training set. 
A batch size of 60 is used, and the model is trained for 32 epochs with a learning rate of 0.01. The Python implementation is shown in listing \ref{lst:training}.

 \begin{lstlisting}[
    language=Python, 
    caption={Model training implementation},
    label={lst:training},
    breaklines=true,
    basicstyle=\ttfamily\small,
    frame=single
]
lr = MultinomialLogisticRegression(train_X.shape[1], train_Y.shape[1])
lr.train(train_X, train_Y, test_X, test_Y, epochs=32, learning_rate=0.01, batch_size=60)
\end{lstlisting}

\subsection*{4.d: Performance evaluation}
The performance of the trained model is evaluated in terms of accuracy by writing to a log file after each epoch. The log file records the 
training and test set log-likelihood and accuracy after each epoch, allowing for analysis of the model's performance over time. The log file is created in a "log" directory, and the results are appended after each epoch of training.
The listing \ref{lst:logging} shows the implementation of the logging functionality.

 \begin{lstlisting}[
    language=Python, 
    caption={Logging implementation for performance evaluation},
    label={lst:logging},
    breaklines=true,
    basicstyle=\ttfamily\small,
    frame=single
]
with open(log_file, "a") as f:
    f.write(f"{epoch+1},{train_ll},{test_ll},{train_acc},{test_acc}\n")
\end{lstlisting}

After executing the training process, the log file can be analyzed to evaluate the model's performance for each epoch. The trends observed provide insights into how well the model generalizes to unseen data and whether it is converging towards a good solution.
Table \ref{tab:model_performance} summarizes the log file data for the last five epochs of training. (A batch size of 60 is used, and the model is trained for 32 epochs with a learning rate of 0.01.)

\begin{table}[h]
    \centering
    \caption{Model Training and Test Performance}
    \label{tab:model_performance}
    \begin{tabular}{c S[table-format=-1.4] S[table-format=-1.4] S[table-format=1.4] S[table-format=1.4]}
        \toprule
        \textbf{Epoch} & {\textbf{Train LL}} & {\textbf{Test LL}} & {\textbf{Train Acc}} & {\textbf{Test Acc}} \\
        \midrule
        28 & -0.3019 & -0.2935 & 0.9163 & 0.9193 \\
        29 & -0.3007 & -0.2926 & 0.9165 & 0.9194 \\
        30 & -0.2995 & -0.2918 & 0.9169 & 0.9195 \\
        31 & -0.2984 & -0.2910 & 0.9174 & 0.9196 \\
        32 & -0.2974 & -0.2902 & 0.9176 & 0.9198 \\
        \bottomrule
    \end{tabular}
\end{table}
The results in the log file indicate that the training and test set log-likelihoods are improving (becoming less negative) over epochs, which corresponds to the model learning to maximise the training log-likelihood. 
Furthermore, it can be observed that both the training and test accuracies are increasing over epochs, indicating that the model is improving its performance on both the training and unseen test data. 
The final test set accuracy of 0.9198 suggests that the model is performing well in classifying the digits in the MNIST dataset.

\subsection*{Question 5: Hyperparameter search}
Optimal hyperparameters for training the model are determined by performing an exploratory search over different batch sizes and numbers of epochs.
The batch size determines the number of samples utilized in one iteration to update the model weights. Small batch sizes provide more frequent updates, whichcan lead to faster convergence, whereas 
larger batch sizes can provide more stable updates but may require more epochs to converge.
The number of epochs defines how many time the entire training dataset is passed through during model training. With too few epochs, underfitting will occur, where the model has insufficient iterations to minimise the log-likelhood. 
With too many epochs, overfitting may occur, where the model learns to fit the training data too closely and fails to generalise to unseen instances of data.  Thus, the optimal number of ephochs is required so as to ensure the underlying distribution of the data is captured effectively.
The batch size values that are tested include 16, 64, 256, and 1024. This follows a logarithmic scale that covers small, medium, and large batch sizes. The epoch counts tested include 10, 50, and 100. 
After training the model with each combination of batch size and epoch count, the training and test accuracies are logged to a CSV file for analysis. The results are summarized in the table \ref{tab:hyperparameter_results} below.

\begin{table}[h]
    \centering
    \caption{Performance Comparison across Batch Sizes and Epochs}
    \label{tab:hyperparameter_results}
    \begin{tabular}{r c S[table-format=1.4] S[table-format=-1.4] S[table-format=1.4]}
        \toprule
        \textbf{Epochs} & \textbf{Batch Size} & {\textbf{Train Acc}} & {\textbf{Test LL}} & {\textbf{Test Acc}} \\
        \midrule
        10  & 16   & 0.9178 & -0.2875 & 0.9194 \\
        50  & 16   & 0.9287 & -0.2680 & 0.9245 \\
        100 & 16   & 0.9322 & -0.2655 & 0.9261 \\
        \midrule
        10  & 64   & 0.9030 & -0.3347 & 0.9091 \\
        50  & 64   & 0.9208 & -0.2823 & 0.9214 \\
        100 & 64   & 0.9256 & -0.2726 & 0.9234 \\
        \midrule
        10  & 256  & 0.8804 & -0.4471 & 0.8880 \\
        50  & 256  & 0.9063 & -0.3239 & 0.9115 \\
        100 & 256  & 0.9144 & -0.2984 & 0.9181 \\
        \midrule
        10  & 1024 & 0.8458 & -0.7277 & 0.8550 \\
        50  & 1024 & 0.8846 & -0.4225 & 0.8917 \\
        100 & 1024 & 0.8969 & -0.3630 & 0.9031 \\
        \bottomrule
    \end{tabular}
\end{table}

Considering the results obtained from the hyperparameter search, it can be observed that the best performance in terms of test accuracy is achieved 
with a batch size of 16 and 100 epochs, yielding a test accuracy of 0.9261. 

It can be observed that smaller batch sizes (16) reach higher accuracies faster than larger batches (1024). This is likely because smaller 
batches provide more frequent updates to the model weights, allowing for faster convergence. It can also be observed that increasing the number of epochs leads to improved performance. (except when overtraining occurs)

\subsection*{Logging}
Don't overwrite files 
Logging throughout training vs logging only at the end of training.


\end{document}
